% Folder 33, batch 03 — archivists' pages 41–60.
% Pass: Opus 5.5 (claude-opus-5-5), 2026-10-03 — first pass, unchecked against the pages by a human.
%
% Handwritten notes on the backs of typescript leaves. The even pages
% (42, 44, 46, 48, 50, 52, 54, 56, 58, 60) carry only typescript unrelated
% to the notes in hand (an English text on sheaves and perfect categories,
% pp. 24-39 of its own numbering; a page on root systems; an EGA IV galley,
% 19.6, with one marginal note of his on graded modules) and are skipped.
% The notes run 41, 43, 45, 47, 49, 51, 53, 55, 57, 59: the cohomology with
% proper supports R_! f_*, base change, global duality for a smooth
% compactifiable morphism, duality for an immersion, relative purity and
% smooth base change as consequences of global duality. The hand is a fast
% draft; most connective prose does not read and is left \ill{}.
\documentclass[11pt,a4paper]{article}
\input{../preamble/grothendieck}

\folder{33}
\batch{3}
\pages{41}{60}
\foldertitle{SGA 1965. MS [Manuscrit] brouillon : notes manuscrites (s.d.).}
\dating{1965-[vers 1971]}
\watermark{Édition de démonstration}

\begin{document}

\page{41}
\note{l'argument court sans doute depuis les pages qui précèdent le lot ; cette page s'ouvre sur une définition sans préambule.}

\emph{Définition des} $\mathbb{R}_! f_*(F^\cdot)$, ($f$ \ill{}).

\begin{tikzcd}
X \arrow[r, hook, "i"] \arrow[d, "f"'] & \overline{X} \arrow[dl, "\overline{f}"] \\
Y &
\end{tikzcd}

\[
\mathbb{R}_! f_* : D^+(X) \longrightarrow D^+(Y), \qquad
\mathcal{H}^i(\mathbb{R}_! f_*) = \mathbb{R}^i_! f_* .
\]

Reformulation de propriétés connues :

(i) $\mathbb{R}_!(gf)_* = (\mathbb{R}_! g_*)(\mathbb{R}_! f_*)$
\uncertain{si} $f$, $g$ \ill{} \uncertain{compactifiables} \ill{}
\uncertain{compatibles}.

(ii) Compatibilité avec extension de la base.

\emph{Remarque.} En fait, si \struck{$f$ \ill{} \ill{}} \add{$Y$ est
\uncertain{quasi-compact}}, on peut définir
\[
\mathbb{R}_! f_* : D(X) \longrightarrow D(Y)
\]
\uncertain{induisant}
\[
D^-(X) \longrightarrow D^-(Y), \quad D^+(X) \to D^+(Y), \quad
D^b(X) \longrightarrow D^b(Y).
\]
Cela résulte du fait que $f_*$, \uncertain{si} $f$ \ill{} \uncertain{de
dimension} \uncertain{bornée}, est de \uncertain{dimension} finie.

\page{43}
Ici, prenons les top. étales sur les préschémas, et les faisceaux de
$A_X$ ($\mathbb{Z}/n\mathbb{Z} = A$), \ill{} \ill{} \ill{} \ill{} de type
\ill{} \ill{} \ill{} \ill{} \uncertain{plats}.

\begin{tikzcd}
X \arrow[d, "f"'] & X' \arrow[l, "h"'] \arrow[d, "{f'}"] \\
Y & Y' \arrow[l, "g"]
\end{tikzcd}

Si $f$ quelconque, on définit, \uncertain{pour} $F^\cdot \in \mathrm{Ob}\, D^+(X)$,
un \uncertain{homomorphisme} \ill{} \uncertain{dans} $D^+(Y')$ :
\[
\text{\struck{$\mathbb{R} f'_*(h^*(F^\cdot)) \longleftarrow$}}
\qquad
\boxed{\,g^*(\mathbb{R} f_*(F^\cdot)) \longrightarrow \mathbb{R} f'_*(h^*(F^\cdot))\,}
\]
\note{le symbole sous $h^*$, dans la formule encadrée, est surchargé ; le premier membre barré est raturé en entier.}

Passant aux faisceaux de coh., cela \uncertain{redonne}
\[
(*_i) \qquad g^* \mathbb{R}^i f_*(F) \longrightarrow \mathbb{R}^i f'_*(h^*(F)).
\]
Lorsque \emph{$f$ est propre}, \ill{} $(*)$ \ill{} \ill{} \ill{}
\uncertain{les} $(*_i)$ \ill{} \ill{}.

\page{45}
\section{§ (9) Le th. de dualité globale : énoncé général pour un morphisme lisse compactifiable}
\note{titre souligné de sa main ; la page ne porte rien d'autre.}

\page{47}
20) \emph{Exemple} \struck{\ill{}} avec un corps $k$ base, alg. clos.
\[
H^i_!(X, F)^\vee \simeq H^{-i}(X, DF).
\]
Si $F$ loc. ct.
\[
H^i_!(X, F)^\vee \simeq H^{2d-i}(X, \underline{\mathrm{Hom}}(F, \mu_{n,X}^{\otimes d})).
\]
\note{l'exposant $2d-i$ est lu d'après la formule semblable plus bas ; ici le trait pourrait aussi se lire $2n-i$.}
On retrouve le th. de finitude de Artin.

Si $S$ quelconque, \ill{} \ill{} \ill{} \ill{} \ill{} loc. constant
(\uncertain{p. ex.} $f$ \uncertain{propre lisse} !)
$R^i_! f_*(F)$ \ill{} \ill{} :
\[
\underline{\mathrm{Hom}}(R^i_! f_*(F), (\mathbb{Z}/n\mathbb{Z})_X) \simeq R^{-i} f_*(D(F))
\]
\note{sous $R^{-i}$, ici et dans la formule suivante, un petit signe surchargé, peut-être un ${}_!$ biffé.}
(ce qui implique que $R^{-i}_! f_*(DF)$ \ill{} loc. constants) et plus
spécialement, si $F$ \uncertain{est} loc. constant :
\[
\underline{\mathrm{Hom}}(R^i_! f_*(F), (\mathbb{Z}/n\mathbb{Z})_X)
\simeq R^{2d-i} f_*(\underline{\mathrm{Hom}}(F, \mu_{n,X}^{\otimes d})).
\]

Donc dans ce cas, \ill{} \uncertain{dualité} \ill{} \ill{}
\add{\uncertain{constructibilité}} \struck{\uncertain{finitude}} des
$R$\struck{\ill{}} \struck{$\underline{\mathrm{Hom}}$} \ill{}
($G = DF$ …).

On peut généraliser \ill{} \uncertain{en supposant} que l'hyp. \ill{}
\uncertain{satisfaite} que pour $i \le i_0$, la \uncertain{conclusion}
\ill{} \uncertain{vaut} pour $j \ge 2d - i_0$ …

\page{49}
\marginal{\uncertain{Le domaine} \ill{} \ill{} \uncertain{étale}}

\section{21. L'homom. de dualité pour une immersion}

Soit $i : Y \to X$ \ill{}, \ill{} \ill{}. Alors $i_!$ est exact (en fait,
identique à $i_*$ si $i$ est une immersion fermée), donc $\mathbb{R} i_!$ est
le \uncertain{foncteur} trivial
\[
\mathbb{R} i_! = i_! : D^+(Y) \longrightarrow D^+(X)
\]
(inutile ici de prendre des \uncertain{résolutions injectives}).

D'autre part, on définit un foncteur
\[
i^! : C(Y) \longrightarrow C(X)
\]
\note{le sens de la flèche est celui de la page ; $i^!$ va de $X$ vers $Y$ d'après la formule qui suit.}
de la façon suivante : on \uncertain{factorise} $i$ en un composé
$Y \xrightarrow{i'} \mathcal{U} \xrightarrow{j} X$, avec $i'$ une
\uncertain{immersion} \emph{fermée}, $j$ une \uncertain{immersion ouverte},
et on pose
\[
i^!(F) = \text{\struck{$\mathbb{R}$}}\; i'^* \underline{\Gamma}_Y(j^*(F)).
\]
Le foncteur ne dépend pas, à iso. unique près, du choix de la factorisation
choisie, comme il \uncertain{résulte} \uncertain{par ex.} \ill{} de la
formule d'adjonction
\[
\mathrm{Hom}_Y(G, i^!(F)) \simeq \mathrm{Hom}_X(i_!(G), F)
\]
qui dit que $i^!$ \struck{\ill{}} est adjoint à droite de $i_!$. Comme
$i_!$ est exact, \ill{} \ill{} \uncertain{injectifs}. Compte tenu de cette
\ill{} \struck{\ill{}} \ill{} \ill{} l'\uncertain{isomorphisme} \ill{} $j_!$
et $j^*$ (\uncertain{fait}
\note{la phrase se poursuit en tête de la p.~51.}

\page{51}
\uncertain{général} pour les topos et \ill{} d'iceux), \ill{} \ill{} \ill{}
\uncertain{cas}.
\[
i_! = (j i')_! = j_! i'_! , \qquad i^! = i'^! j^! = i'^! j^* .
\]
On \uncertain{revient} à définir \struck{\ill{}} \struck{\ill{} \ill{}}
$i^!$ comme foncteur adjoint de $i_!$, \ill{} \uncertain{au} cas où $i$
est une imm. \emph{fermée}, \ill{} \ill{} \ill{} \ill{} \uncertain{fait}
\uncertain{général} \ill{} pour les topos et \ill{} d'iceux,
cf Exp IV.
\struck{\ill{}} Comme $i_!$ est exact, \ill{} \ill{} \ill{}
\uncertain{adjoint} à \uncertain{droite} \ill{} $i^!$, \ill{} \ill{}
\uncertain{injectifs}, donc \ill{} \ill{} \emph{\uncertain{injectif}},
\struck{\ill{}} \ill{} \ill{}
\[
\mathbb{R} i^! : D^+(X) \longrightarrow D^+(Y).
\]
\ill{} \ill{} \ill{} \ill{} \ill{}, \ill{} \ill{} que le
\uncertain{foncteur} \emph{\ill{}} \ill{} $\mathcal{U}$
(\struck{\ill{}} \ill{} cat. \uncertain{dérivées}). \struck{\ill{}}

\struck{\ill{}}
\note{un passage de trois lignes, encadré et biffé de traits ondulés, se termine sur $\mathbb{R}\underline{\Gamma}_Y$ ; on n'en lit pas davantage.}

Notons, \ill{} $i$ est une imm. \emph{fermée},
\[
(\mathbb{R}_! i_*)(\mathbb{R} i^!) \simeq \mathbb{R} \underline{\Gamma}_Y
\]
\note{le facteur de gauche est surchargé : un ${}^*$ récrit en ${}_*$, entre parenthèses ajoutées.}
\ill{} \ill{}, \ill{} \ill{} \ill{} \ill{}, d'où \ill{} si
$F \in C(X)$, \ill{} \ill{}
\[
\underline{H}^q(\mathbb{R} i^!(F)) \simeq i^*(\underline{H}^q_Y(F)).
\]

\page{53}
Du fait que $i^!$ est l'adjoint à droite \struck{commutant aux produits}
\uncertain{d'un} foncteur exact \struck{\ill{}} ($i_!$) \ill{} la
\emph{formule} de \emph{dualité pour les immersions} :
\[
\boxed{\,\mathrm{Hom}(G^\cdot ; \mathbb{R} i^!(F^\cdot)) \simeq
\mathrm{Hom}(\mathbb{R}_! i_*(G^\cdot), F^\cdot)\,}
\]
(isom. bifonctoriel).

\begin{align*}
\mathrm{Hom}(G^\cdot ; \mathbb{R} i^!(F^\cdot))
&= \mathrm{Hom}(G^\cdot ; i^!(C^\cdot(F^\cdot))) \\
&= H^0(\mathrm{Hom}^\bullet(G^\cdot ; i^!(C^\cdot(F^\cdot))))
 = H^0(\mathrm{Hom}(i_!(G^\cdot), C^\cdot(F^\cdot))) \\
&\text{\struck{$= \mathrm{Hom}(i_!(G^\cdot), \ldots)$}} \\
&= \mathrm{Hom}(\mathbb{R} i_!(G^\cdot), C^\cdot(F^\cdot))
\end{align*}
\note{le $H^0$ de la deuxième égalité de la deuxième ligne est ajouté au-dessus de « Hom ».}

Quand on explicite la flèche $\to$, on trouve que l'\ill{}
\uncertain{morphisme} \ill{} un \uncertain{homom.} can.
\[
\mathrm{Tr}_i(F^\cdot) : \mathbb{R}_! i_* \mathbb{R} i^!(F^\cdot) \longrightarrow F^\cdot
\]
de la \uncertain{même} façon que dans le cas \uncertain{du} th. de dualité
globale pour un morphisme lisse.

\page{55}
\section{22. Th. de pureté relative comme conséquence du th. de dualité globale}
\note{dans le titre, un mot biffé et encadré est remplacé par « relative », de sa main.}

\begin{tikzcd}
Y \arrow[r, hook, "i"] \arrow[dr, "g"'] & X \arrow[d, "f"] \\
& S
\end{tikzcd}

Supposons que [\ill{} $i$ une \uncertain{immersion} \uncertain{fermée}]
\ill{} $g$ soit \struck{\ill{}} \struck{\ill{}} \ill{}, alors la formule
de \uncertain{dualité} \ill{} \ill{} :
\[
\mathbb{R}_! f_*\bigl(\mathbb{R}\underline{\mathrm{Hom}}(A_Y, \mathbb{R}^! f(G^\cdot))\bigr)
\simeq \mathbb{R}\underline{\mathrm{Hom}}(\mathbb{R}_! g(A_Y), G^\cdot)
\]
\note{une accolade sous le terme $\mathbb{R}\underline{\mathrm{Hom}}(A_Y, \mathbb{R}^! f(G^\cdot))$ le renvoie à $\mathbb{R} i_* \mathbb{R} i^!(\mathbb{R}^! f(G^\cdot))$, écrit dessous.}
s'écrit
\[
g_*\bigl(\mathbb{R} i^!(\mathbb{R}^! f(G^\cdot))\bigr) \xrightarrow{\ \sim\ } G^\cdot ,
\]
\ill{} \ill{} $G^\cdot$ \ill{} \ill{} \ill{} \ill{}
\[
\begin{cases}
\underline{H}^i_Y(f^*(G)) = 0 & \text{si } i \neq 2d \\
\underline{H}^{2d}_Y(f^*(G)) = g^*(G) \otimes T^{-1}_{Y/X} &
\end{cases}
\]
\struck{\ill{}} \ill{}.

\struck{\ill{} \ill{}}, \struck{\ill{} \ill{}} \ill{} maintenant
$Y \subset X$ \ill{} \ill{} \ill{} : \ill{} \ill{} que $Y$, $X$
\ill{} \ill{} \ill{} \ill{} $d$ \ill{} \ill{} \ill{} \ill{} \ill{}, $X$
quelconque, \ill{} \struck{\ill{}} \ill{} $\Rightarrow$ \ill{}
\ill{} \ill{} \ill{} \ill{} \ill{} \uncertain{flèches} \ill{}
\ill{} \ill{} \ill{} \ill{} \ill{} $X$ \ill{} \ill{} \ill{}
\uncertain{question} \ill{}, \ill{} \ill{}, \ill{} \ill{}
\uncertain{considérons} \ill{} \ill{} :

\begin{tikzcd}
Y \arrow[r, hook] \arrow[d] & X \arrow[d] \\
Y' \arrow[r] \arrow[dr, "{g'}"'] & X' = S[t_1, \ldots, t_d] \arrow[d, "{f'}"] \\
& S
\end{tikzcd}

\struck{et \ill{}} avec $X \to X'$ \emph{lisse},
\note{la phrase se poursuit en tête de la p.~57.}

\page{57}
(N.B. $Y'$ est la « section nulle » de $S[t_1, \ldots, t_d]$) et compte tenu
du th. de \emph{changement de base lisse}, on est ramené au cas précédent.

\emph{Th. de changement de base lisse} comme conséquence du th. de dualité
globale (d'après Verdier). \ill{} \uncertain{facilement}, \ill{}
\uncertain{la} \uncertain{formule} (compte tenu du changement de base
\uncertain{propre}) \ill{} \ill{} \ill{} \uncertain{d'un} \uncertain{ouvert}
$\mathcal{U} \to X$, \ill{} et \ill{} \ill{} \ill{} \ill{} $X' \to X$, \ill{}
\ill{} \ill{} \ill{} ; \ill{} \ill{} \ill{}, \ill{} \ill{}
$Y = X - \mathcal{U}$)
\note{dans la marge gauche, à hauteur de ces lignes, une flèche verticale griffonnée, sans lettre lisible.}

\begin{tikzcd}
Y \arrow[d, "i"'] & Y' \arrow[l, "g"'] \arrow[d, "{i'}"] \\
X & X' \arrow[l, "f"]
\end{tikzcd}

\[
g^*(\mathbb{R} i^!) \longrightarrow (\mathbb{R} i'^!) f^*
\]
$T_{Y'/Y}$ \ill{}. \ill{} \ill{} \ill{} \ill{}, \ill{} \ill{} \ill{} \ill{}.

\ill{} \ill{} \ill{} \uncertain{naturel}
\[
(\mathbb{R}^! g^*)(\mathbb{R} i^!) \longrightarrow (\mathbb{R} i'^!)(\mathbb{R}^! f)
\]
\ill{} \ill{} \ill{} un \uncertain{isom}. \ill{} \ill{} \uncertain{dualité}
\ill{} \uncertain{essentiellement} que pour $f$, $g$ \ill{}, \ill{} \ill{}
\ill{} \ill{} « \uncertain{adjoints} » de
$(\mathbb{R}_! i_*)(\mathbb{R}_! g)$ et de
$(\mathbb{R}_! f)(\mathbb{R}_! i')$, \ill{} \uncertain{sont}
\uncertain{évidemment} \uncertain{isomorphes}. Il faut évidemment
\uncertain{prouver} que l'\struck{\ill{}} \uncertain{iso.} \struck{\ill{}}
précédent est bien l'« \uncertain{adjoint} » de l'\ill{} $(*)$, ce que je
\uncertain{n'ai} pas \uncertain{fait}. Je \ill{} \ill{} \ill{} \ill{}
\uncertain{de} \ill{}, car il faudrait \uncertain{savoir}, lorsque \ill{}
\ill{} \ill{} \uncertain{argument} $F \in \mathrm{Ob}\, C(X)$
\note{la phrase se poursuit en tête de la p.~59.}

\page{59}
\note{en tête, un « 4 » isolé.}
\struck{\ill{}} \ill{} \ill{} \ill{} \ill{} \ill{}
\note{cette première ligne est cerclée et renvoyée par un trait à la suivante.}
\ill{} les deux \ill{} \ill{} \ill{} \struck{\ill{}} \ill{} $F$, \ill{}
\uncertain{bien} \ill{} \uncertain{bornés}, \ill{} \uncertain{général}),
ce qui me semble \uncertain{nullement} clair. Mais \ill{} \ill{} \ill{}
\ill{} \ill{}, \ill{} \ill{} \ill{} \ill{} \ill{} \ill{} \ill{}
$X$ \ill{} \ill{} \uncertain{type} \uncertain{fini} \ill{} $\mathbb{Z}$.
\struck{et \ill{} \ill{} \ill{} \ill{} $S[t_1, \ldots, t_n]$}.

Alors le th. de \ill{} \ill{} \ill{} \ill{} \ill{} \ill{} \ill{} \ill{}
$\mathbb{R} f_*$, \ill{} \ill{} \ill{} \ill{} \ill{} \ill{} \ill{}

\struck{\ill{}}
\note{trois lignes, un mot et une formule surchargés, sont biffées d'un trait ondulé, avec des points d'interrogation en fin de ligne.}

Ref exposé d'Artin.
\note{la note s'arrête ici, sur un renvoi ; la suite de l'argument, s'il y en a une, est au-delà du lot.}

\end{document}
